% \chapter{Methods}


\chapter{\name{}}

\label{sec:approach}

\begin{comment}
The most naive approach to the above synthesis problem is, to enumerate all possible symbolic actions vectors, fill in symbolic values with all possible parameters within a given range, and execute all these attack vectors in a forked simulation environment to check if they actually yield positive profit. However, this solution is infeasible in many aspects. (1) Large range of parameters. Flash loan attacks usually utilize large amounts of assets which indicates large range of parameters. (2) Slow execution. It is notoriously slow to fork a block and execute our transactions in a locally simulated environment.~\cite{kim2021off}


Another approach is to manually extract closed-form expressions of state transition functions, construct the optimization framework for an attack vector and solve the optimization problem mathematically, as shown in prior work~\cite{qin2021attacking,cao2021flashot}. However, this approach still require (1) an expert knowledge about DeFi protocols and smart contracts (2) is not applicable to DeFi attacks involving smart contracts that are not open-sourced, such as the case for the Harvest USDC attack described in Section~\ref{sec:example} (3) is not applicable which it is not possible to extract the closed-form expression of a DeFi protocol endpoint. For example, the decentralized exchange Curve ~\cite{egorov2019stableswap} uses an iterative method to calculate stablecoins' prices while maintaining its StableSwap invariant. It is impossible to extract a closed form expression of stablecoins' prices for a pool with $\geq 3$ tokens as it requires solving a three-variable cubic equation. 


To address the above concerns, we present a framework for attack vectors synthesis which first searches for all  symbolic actions vector that can lead to feasible attacks, approximates state transition functions, then constructs the optimization framework to automatically find for all possible attack vectors.
\end{comment} 


%\subsection{Symbolic Attack Vector Synthesis}\label{sec:symbolicSynthesis}

\begin{comment}
\begin{algorithm}[!h]

\caption{Attack vectors synthesis procedure.}
\alglabel{alg:enu}
\SetAlgoLined
\DontPrintSemicolon
\small
% -------------------------------------------------------
\SetKwInOut{Input}{Input}
\SetKwInOut{Output}{Output}
\SetKwInOut{Procedure}{Procedure}
% -------------------------------------------------------
\SetKwData{w}{\textit{worklist}}

\SetKwData{wp}{\textit{wlist}}
\SetKwData{al}{\textit{answerlist}}
\SetKwData{e}{\textit{e}}
\SetKwData{E}{\textit{E}}
\SetKwData{p}{\textit{p}}
\SetKwData{pp}{\textit{p}}
\SetKwData{s}{\textit{s}}
\SetKwData{d}{\textit{d}}
% -------------------------------------------------------
\Procedure {Synthesize($\mathbf{A}$, $len$, $\mathcal{P}$, $q$, $it$)}
\Input{
   a set of actions $\actionset$, a maximum length of an actor vector $len$, a target profit function $\mathcal{P}$, a blockchain state $q$
}
\Output{
  sequences of actions with corresponding parameters which yield the best target profit
}
% -------------------------------------------------------
\For {a in $\mathbf{A}$} { \label{algo1:line0}
    datapoints[a] \assign dataCollect($q$, a)\\ \label{algo1:line1}
} 
it \assign 0\\
\For {it < It}  { \label{algo1:line2}

    $\mathbf{A}'$ \assign approximate($\mathbf{A}$,datapoints) \\ \label{algo1:line3}
    \wp \assign actionsVectors($\mathbf{A'}$, $len$) \\ \label{algo1:line4}

    \For {\pp in \wp} { \label{algo1:line5}

        \eIf{IsFeasible(\pp)} { \label{algo1:line6}
            ($\pp^\star$, profit) \assign optimize(\pp, $\mathcal{P}$) \\ \label{algo1:line7}
            \eIf{queryOracle($q$, $\pp^\star$, profit)} { \label{algo1:line8}
                \al.add(($\pp^\star$,profit)) \\ \label{algo1:line9}
            }
            {
                datapoints \assign datapoints $\cup$ CEGDC($\pp^\star$, $q$)\\ \label{algo1:line10}
            }            
        } { continue }    
    }
    it \assign it + 1 \label{algo1:line11}
}
\Return \al
\end{algorithm}
\end{comment}

In this chapter, we present {\name} framework for attack vectors synthesis. {\name} first searches all  symbolic actions vector that can lead to feasible attacks and approximates state transition functions. {\name} then constructs the optimization framework to automatically find for all possible attack vectors. 

In Algorithm~\ref{alg:enu}, we give the synthesis procedure of {\name}.  The procedure first collects initial data points to approximate the actions in $\actionset$ (line~\ref{algo1:line1}) where we use as a starting blockchain state the state $\q$. Then, using the sub-procedure \textsc{Approximate} we generate the approximation of the actions in $\actionset$ using the collected data points (line~\ref{algo1:line3}). We use the sub-procedure \textsc{ActionsVectors} to generate all possible symbolic actions vectors of length less than $\len$ (line~\ref{algo1:line4}). 
We iterate over the generated actions vectors where use some refinement heuristics implemented in the sub-procedure \textsc{IsFeasible} to prune actions vectors (line~\ref{algo1:line6}). For instance, an actions vector containing two adjacent actions invoking the same method that swaps token X to token Y can be pruned to an actions vector where the two adjacent actions are merged. Afterwards, we use the optimization sub-procedure \textsc{Optimize} (line~\ref{algo1:line7}) to find the suitable concrete values to pass as input parameters to the methods in the actions vector that satisfy the constraints encoded in the objective function $\mathcal{P}$. In the optimization procedure we use for each method its approximated version. We then validate whether the attack vector generated by the optimizer does indeed generate the profit. To do this we use the sub-procedure \textsc{QueryOracle} to execute the actions vector with the optimizer generated input parameters on the actual smart contracts on the blockchain. If the query is successful, i.e., the actual profit closely matches the profit found by the optimizer, then we add the attack vector to the list of discovered attacks. Otherwise, if the query is not successful, we consider the attack vector to be a counterexample, and we use it to generate new data points to refine the approximation in the subsequent iterations, the sub-procedure \textsc{CEGDC} (line~\ref{algo1:line10}). We repeat the process until the number of iterations reaches a threshold $\iter$ (line~\ref{algo1:line2}).

%the worklist to  $\epsilon$. The loop at lines 3 to line 22 dequeues an actions vector from the worklist(line 4) and expand it using the set of actions $\mathbf{A}$ (line 5). For an actions vector \textit{p'} in the expansion result \textit{worklist'}, we process \textit{p'} based on whether it is partial or symbolic. If it is partial (line 7), we add it to the \textit{worklist}. If it is symbolic (line 9), we invoke a procedure IsFeasible (described in Section~\ref{sec:isfeasible}) to check if \textit{p'} will lead to a feasible attack. If it is, we automatically construct an optimization framework (described in Section~\ref{sec:opframe}) and use an optimizer to find an attack vector $\attackvector^\star$ that generates positive profit (line 11), i.e., answer candidates. Then we query an oracle to execute $\attackvector^\star$ (line 12), and get $\attackvector^+$ (set of concrete attack vectors whose execution results shows positive profit) and $\attackvector^-$(a set of concrete attack vectors whose estimated profit is different from execution results). $\attackvector^+$ is what we want, and we directly add it to the \textit{answerlist}(line 13). For each concrete attack vector in $\attackvector^-$, we invoke a procedure CEGDC(described in Section~\ref{sec:cegdc}) to get data points which need to be added to refine our transition functions' approximations(line 15). 

\setlength{\textfloatsep}{10pt}
\begin{algorithm}[h!]
  \caption{Attack vectors synthesis procedure}
  \label{alg:enu}
  
  \textbf{Inputs: } a set of actions $\actionset$, the maximum length of an actor vector $\len$, an objective function $\mathcal{P}$, a blockchain state $\q$, and a threshold number of iterations $\iter$.
  
  \textbf{Outputs: } sequences of actions with corresponding parameters that yield positive profits.
  \begin{algorithmic}[1]
  \Procedure{Synthesize}{$\actionset$, $\len$, $\mathcal{P}$, $\q$, $\iter$}
  \State \ \ \textbf{for each}\ $\act \in \actionset$ \label{algo1:line0}
  \State \ \ \ \ $\mathsf{datapoints}[\act] \leftarrow \textsc{DataCollect}(\q, \act);$ \label{algo1:line1}
  \State \ \ \textbf{for each}\ $i \in [0;\ \iter]$ \label{algo1:line2}
  \State \ \ \ \ $\approxact \leftarrow \textsc{Approximate}(\actionset,\mathsf{datapoints});$ \label{algo1:line3}
  \State \ \ \ \ $\wkp \leftarrow \textsc{ActionsVectors}(\approxact, \len);$ \label{algo1:line4}
  \State \ \ \ \ \textbf{for each}\ $\pp \in \wkp$ \label{algo1:line5}
  \State \ \ \ \ \ \ \textbf{if}\ $\textsc{IsFeasible}(\pp)$ \label{algo1:line6}
  \State \ \ \ \ \ \ \ \ $(\pp^\star, \textsf{profit}) \leftarrow \textsc{Optimize}(\pp, \mathcal{P})$; \label{algo1:line7}
  \State \ \ \ \ \ \ \ \ \textbf{if}\ $\textsc{QueryOracle}(\q, \pp^\star, \textsf{profit})$ \label{algo1:line8}
  \State \ \ \ \ \ \ \ \ \ \ $\al.\textsf{add}(\pp^\star,\textsf{profit})$; \label{algo1:line9}
  \State \ \ \ \ \ \ \ \ \textbf{else} 
  \State \ \ \ \ \ \ \ \ \ \ $\mathsf{datapoints} \assign \mathsf{datapoints} \cup \textsc{CEGDC}(\pp^\star, \q)$; \label{algo1:line10}
  \State \ \ \textbf{return}\ $\al$;
  \EndProcedure
  \end{algorithmic}
\end{algorithm}


    
    

\section{Pruning Symbolic Actions Vectors} \label{sec:isfeasible} 
The sub-procedure \textsc{IsFeasible} implements some heuristics to prune undesired symbolic actions vectors. We present here some heuristics, and elaborate more on others in Section~\ref{sec:implementation}. 

\noindent \textbf{Heuristic 1: no duplicate adjacent actions.} Two successive calls the same method in a DeFi smart contract twice are usually equivalent to a single call to the same method with different parameters. Thus, we discard actions vectors that contain two successive calls to the same method.

\noindent \textbf{Heuristic 2: limited usage of a single action.} Using the observation that for attack vectors that do not contain repetitions, e.g., loops, one method is called limited number times. DeFi attacks usually involve methods from different smart contracts. Thus, we fix a maximum number of calls to a single method an attack can contain and discard actions vectors that do not satisfy this criterion, e.g., an actions vector of length $4$ cannot contain more than two calls to the same method.  




%However, it is impossible to know in advance the number of data points needed to well approximate a state transition functions; thus, we design a counterexample driven data collection(CEGDC) loop to argument data points when an approximation error is identified. 

%The high level idea of approximating state transition functions is, with enough data points of $((\pstate, p), (\pstate', Q_{in}, Q_{out}) )$, we can infer an approximation of $f_1, f_2, f_3$ without the need to understand the underlying mathematical models of related DeFi protocols. Given a set of data points $((\pstate, p), (\pstate', Q_{in}, Q_{out}) )$, we aim to find approximating functions $f_1$, $f_2$, $f_3$, such that $\pstate' \approx f_1(\pstate, p)$, $Q_{in} \approx f_2(\pstate, p)$, and $Q_{out} \approx f_3(\pstate, p)$.


%An LTS is \emph{deterministic} if for any state $s$ and sequence $\tau\in \Sigma^*$, there is at most one state $s'$ such that $s\xrightarrow{\tau}s'$. 

%\noindent \textbf{Key Observation:} For any state transition function of a DeFi protocol endpoint, $\transit(p)$:  $(\pstate,\balance)\xrightarrow[]{\transit(p)}(\pstate', \balance')$, given the same protocol prestates $\pstate$ and input parameter $p$, the effect of calling this endpoint will always be the same, in other words, $\pstate'$ and the difference of $\balance'$ and $\balance$ will be the same, unless the transaction is reverted. 



%\begin{itemize}
%    \item protocol poststates $\pstate'$.  Calling a DeFi endpoint will likely change the states of one or more DeFi protocols. $\pstate' = f_1(\pstate, p)$
%    \item token consumed $Q_{in}$. Some DeFi endpoints transfer tokens from the endpoint caller, it can be written as $Q_{in} = f_2(\pstate, p)$ 
%    \item token got $Q_{out}$. Some DeFi endpoints transfer tokens to the endpoint caller, it can be written as $Q_{out} = f_3(\pstate, p)$
%\end{itemize}
%
%The high level idea of approximating state transition functions is, with enough data points of $((\pstate, p), (\pstate', Q_{in}, Q_{out}) )$, we can infer an approximation of $f_1, f_2, f_3$ without the need to understand the underlying mathematical models of related DeFi protocols. Given a set of data points $((\pstate, p), (\pstate', Q_{in}, Q_{out}) )$, we aim to find approximating functions $f_1$, $f_2$, $f_3$, such that $\pstate' \approx f_1(\pstate, p)$, $Q_{in} \approx f_2(\pstate, p)$, and $Q_{out} \approx f_3(\pstate, p)$.

%We use two different numerical methods to infer the approximation -- polynomial approximation and multivariate interpolation (described in Section ~\ref{sec:statetransitimple}). 

%However, it is impossible to know in advance the number of data points needed to well approximate a state transition functions; thus, we design a counterexample driven data collection(CEGDC) loop to argument data points when an approximation error is identified. 

\section{Transition Functions Approximation} \label{sec:approximate}
We now describe the sub-procedure \textsc{Approximate} for approximating the behaviors of actions. 
Note that smart contracts are deterministic since nodes in the blockchain relies on deterministic outcomes for any sequence of transactions to reach consensus. Thus, the state reached by contracts when executing a sequence of actions is unique. Therefore, as the effect of an action is solely dependent on the current state $\q$ and input parameters $\textsf{params}$, it is possible to write every element of the post-state as a function of the pre-state $\q$ and the parameters $\textsf{paras}$, i.e., $\q' = \textsf{f}(\q, \textsf{paras})$. In \textsc{Approximate}, we use data-driven approximation to estimate the function $\textsf{f}$ using collected data points that consist of the tuples $(\q',\textsf{paras},\q)$. To carry the estimation we use two different numerical methods, i.e., polynomial approximation and multivariate interpolation. 

%\footnote{note that $f$ can be a vector of functions where each function an element map the state q given the input parameters to an element of the state $q'$.}.

\subsection{Polynomial} 
We use standard polynomial features extraction to generate feature matrices from data points~\cite{statics1}. 
The feature matrices consist of all polynomial combinations of the features with degree less than or equal to some constant $n$. Then, we use linear regression to find polynomials' coefficients which minimize the residual sum of squares between the observed post-states in the data points, and the post-states predicted by the linear approximation~\cite{statics2}. We apply the obtained polynomials to all pre-states and input parameters and distinguish the number of miss-approximated data points which we use to choose the polynomial which gives the fewest miss-approximated data points. %However, not every multivariate function can be approximated using multivariate polynomials. Thus, multivariate interpolation method is also adopted as an alternative to multivariate polynomial approximation method.
%In practice, for any state transition function $f$, we tried $n = 1, 2, ..., 6$, and choose the polynomial which gives us fewest misapproximated points as our desired approximating function. 

\subsection{Interpolation} 
We also use multivariate interpolation method~\cite{interpolation} since not every multivariate function can be approximated using multivariate polynomials. In particular, we use nearest-neighbor interpolation to build the nearest $N$ dimensional interpolator from data points tuples. Afterwards, to return an estimated output for an input value, the interpolator returns the nearest neighbor of the input value. 


\section{Optimization} \label{sec:opframe} 
Given a symbolic actions vector and the approximated methods called by those actions, the sub-procedure \textsc{Optimize} constructs an optimization problem to find optimal values for the parameters for the actions. Recall that given a blockchain state $\q$ and an address $\adr$, the actions vector $\symbolic$ transforms $\q$ to another state $\q'$. The objective function in the optimization problem targets to increase the tokens values in the balance of the address $\adr$, i.e., $y = \balance(\q',\adr) -  \balance(\q,\adr)$. The optimization problem is accompanied by constraints on the symbolic values to be inferred. For instance, the balance of a token must always be non-negative, i.e., the adversary and the smart contracts cannot use more tokens than what they have in their balances, otherwise the transaction reverts. %Also, for a symbolic parameter $p_i$ that we want to infer a value for, we fix an upper bound value $u_i$. 
In the following, we give an example of the optimization problem.
\vspace{-1mm}
\begin{align*}  
\mathcal{P}: \ \ \ & \begin{cases} 
    \max_{p_0, p_1, ..., p_n}\ y = \balance(q',\adr) -  \balance(q,\adr) \\ 
   \textrm{subject to: } \forall\ \mathbf{t} \in \mathbf{T}, \adr' \in \mathbf{A}.\ \tokenmap(q',\adr',\mathbf{t})  \geq 0   
    \end{cases}       
\end{align*}  

\begin{comment}
\begin{equation}
\begin{split}
&\max_{p_0, p_1, ..., p_n}\ y = \balance(q',\adr) -  \balance(q,\adr) \\ 
 \textrm{subject to: } &
    \begin{cases}
      0 < p_i < u_i, i = 0, 1, ..., n \\
      \forall\ \mathbf{t} \in \mathbf{T}, \adr' \in \mathbf{A}.\ \tokenmap(q',\adr',\mathbf{t})  \geq 0   
    \end{cases}       
\end{split}
\end{equation}
\end{comment}


\section{Counterexample Guided Data Collection (CEGDC)} \label{sec:cegdc}
The optimization sub-procedure might explore parts of the states space not explored during the initial data points collection. This might challenge the accuracy of the approximations and result in mismatch between the estimated  and the actual values. Thus, it is necessary to collect new data points based on the counterexamples that show the mismatch between the estimated and the actual values, to refine the approximations. 
This allows to improve the accuracy of approximation and eliminate the mismatches. 
Therefore, we propose counterexample guided data collection (CEGDC), inspired of counterexample guided abstraction refinement~\cite{clarke2000counterexample}, to refine approximations when mismatches are identified. 

We use $\concrete$ to denote the attack vector s.t. $\q \xrightarrow{\concrete} \q'$. 
$\q'_{e}$ and $\q'_{a}$ denote the estimated value for the state $\q'$ found by the optimizer and the actual value obtained when executing $\concrete$ on the actual protocol on the blockchain in the state $\q$, respectively. 
$\profit_e(\concrete) = \balance(\q'_{e},\adr) -  \balance(\q,\adr)$ and $\profit_a(\concrete) = \balance(\q'_{a},\adr) -  \balance(\q,\adr)$ denote the estimated profit and actual profit, respectively. 

\begin{definition} \label{counterexample}
A counterexample is an attack vector $\concrete$ whose estimated profit $\profit_e(\concrete)$ is different from its actual profit $\profit_a(\concrete)$. Formally, $|\profit_e(\concrete) - \profit_a(\concrete)| \geq \varepsilon \cdot (|\profit_e(\concrete)| + |\profit_a(\concrete)|) $, where $\varepsilon$ is a small constant representing accuracy tolerance.
\end{definition}


% $\concrete[:-k]$ denotes the vector constituted of the elements $\concrete[0]$ $\cdots$ $\concrete[len(\concrete)-k]$
\begin{comment}
\begin{algorithm}[!h]

    \caption{Counterexample guided data collection.}
    \label{alg:cegdc}
    \SetAlgoLined
    \DontPrintSemicolon
    \small
    % -------------------------------------------------------
    \SetKwInOut{Input}{Input}
    \SetKwInOut{Output}{Output}
    \SetKwInOut{Procedure}{Procedure}
    % -------------------------------------------------------
    \SetKwData{w}{\textit{worklist}}
    \SetKwData{wp}{\textit{worklist'}}
    \SetKwData{al}{\textit{answerlist}}
    \SetKwData{e}{\textit{e}}
    \SetKwData{E}{\textit{E}}
    \SetKwData{p}{\textit{p}}
    \SetKwData{pp}{\textit{p'}}
    \SetKwData{s}{\textit{s}}
    \SetKwData{D}{\textit{datapoints}}

    % -------------------------------------------------------
    \Procedure {CEGDC($\concrete$,$\q$)}
    \Input{
    A counterexample $\concrete$ and a blockchain state $\q$.
    }
    \Output{
        datapoints
    }
    % -------------------------------------------------------
    datapoints \assign [] \\
    %$\p'$ \assign concretize($\p$,$\concrete$) \\ \label{algo2:line0}
    \For{$k = len(\concrete)$ to $1$} { \label{algo2:line1}
        $q'_e$ \assign estimate($\q$, $\concrete$, k)     \\ \label{algo2:line2}
        $q'_a$ \assign execute($\q$, $\concrete$, k)     \\ \label{algo2:line3}
        %$((\pstate, p), (\pstate', Q_{in}, Q_{out}) )$ = dataCollect($\concrete[len(\concrete)-k]$)  \\

        %$((\bar{\pstate}, \bar{p}), (\bar{\pstate}', \bar{Q_{in}}, \bar{Q_{out}}) )$ = simulate($\concrete[:-k + 1]$) \\
        
        \eIf{isAccurate($q'_{e}, q'_{a}$)} { \label{algo2:line4}
            
            returns  datapoints \\ \label{algo2:line5}
        } 
        { 
            (a,paras) \assign $\concrete[k]$ \\ \label{algo2:line7}
            datapoints[a] \assign (q,paras,$q'_{a}$)  \\ \label{algo2:line8}
        }
    } 
\end{algorithm}
\end{comment}

\setlength{\textfloatsep}{10pt}
\begin{algorithm}[h!]
  \caption{Counterexample guided data collection procedure. It takes a counterexample $\concrete$ and a  state $\q$, and returns datapoints. $i \in [n; 1]$ means that in the first iteration $i = n > 0$. }\label{alg:cegdc}
  \begin{algorithmic}[1]
  \Procedure{CEGDC}{$\concrete$, $\q$}
  \State \ \ $\mathsf{datapoints} \leftarrow [\ ]$;
  \State \ \ \textbf{for each}\ $i \in [len(\concrete);\ 1]$; \label{algo2:line1}
  \State \ \ \ \ $\q'_e \leftarrow \textsc{Estimate}(\q, \concrete, k);$ \label{algo2:line2}
  \State \ \ \ \ $\q'_a \leftarrow \textsc{Execute}(\q, \concrete, k);$  \label{algo2:line3}
  \State \ \ \ \ \textbf{if}\ $\textsc{IsAccurate}(\q'_{e}, \q'_{a})$  \label{algo2:line4}
  \State \ \ \ \ \ \ \textbf{returns}\ $\mathsf{datapoints}$ ; \label{algo2:line5}       
  \State \ \ \ \ \textbf{else}
  \State \ \ \ \ \ \ $(\act,\textsf{paras}) \leftarrow \concrete[k]$; \label{algo2:line7}
  \State \ \ \ \ \ \ $\mathsf{datapoints}[\act] \leftarrow (\q,\textsf{paras},\q'_{a})$; \label{algo2:line8}
  \State \ \ \textbf{return}\ $\mathsf{datapoints}$;
  \EndProcedure
  \end{algorithmic}
\end{algorithm}

In Algorithm~\ref{alg:cegdc}, we present the sub-procedure \textsc{CEGDC} for collecting new data points from a counterexample. \textsc{CEGDC} takes as inputs a counterexample $\concrete$ which is known to have an inaccurate profit estimation, and a blockchain state. The for loop at line~\ref{algo2:line1} is used to locate approximation errors backward from the last action to the first action and collect new data points accordingly. In a loop iteration $k$, we check if the estimated methods of the action at the index $k$ of $\concrete$ are accurate. First, we compute the estimated state $\q'_e$ reached by executing $\concrete$ until reaching the action indexed $k$ (line~\ref{algo1:line2}) using the approximated of transition functions. 
Second, we compute the actual state $\q'_e$ reached by executing $\concrete$ until reaching the action indexed $k$ (line~\ref{algo1:line3}) on the actual smart contracts on the blockchain. 
Then, we compare the estimated and actual execution results (line~\ref{algo1:line4}). 
If the estimation is accurate, this indicates that the transition functions of the action at the index $k$ of $\concrete$ and its predecessors are accurate; so the procedure breaks the loop and returns the data points computed in the previous iterations (line~\ref{algo2:line5}). Otherwise, it indicates inaccurate transition functions of this action or/and its predecessors. Thus, we add a new data point associated with the action at the index $k$ of $\concrete$ (lines~\ref{algo2:line7} and \ref{algo2:line8}) and proceed to the next iteration of the loop to explore the action predecessors.









